\section*{Chapitre \ref{ch:g12:deriv} --- Dérivation et convexité}

\begin{solution}{exo:g12:deriv:1}
$f'(x) = 3x^2 - 5$.

Règle du quotient~:
$g'(x) = \dfrac{(x^2+1) - x\cdot 2x}{(x^2+1)^2} = \dfrac{1 - x^2}{(x^2+1)^2}$.

Règle de la chaîne avec $u = 2x+1$~: $h'(x) = 5 \cdot 2 \cdot (2x+1)^4 = 10(2x+1)^4$.

Règle de la chaîne avec $u = x^2 + 1$~:
$k'(x) = \dfrac{2x}{2\sqrt{x^2+1}} = \dfrac{x}{\sqrt{x^2+1}}$.
\end{solution}

\begin{solution}{exo:g12:deriv:2}
$f'(x) = 2x - 3$, donc $f(2) = -1$ et $f'(2) = 1$~: la \omterm{def:g12:deriv:tangent}{tangente} est
$y = -1 + (x - 2) = x - 3$. L'écart est
\[
f(x) - (x - 3) = x^2 - 4x + 4 = (x-2)^2 \geq 0 ,
\]
donc la courbe est au-dessus de la \omterm{def:g12:deriv:tangent}{tangente} partout, ne la touchant qu'en
$x = 2$ (comme attendu~: $f$ est \omterm{def:g12:deriv:convex}{convexe}).
\end{solution}

\begin{solution}{exo:g12:deriv:3}
Ce sont tous deux des taux d'accroissement. Avec $f(x) = (1+x)^{100}$ en
$0$~: $f'(x) = 100(1+x)^{99}$, donc la limite est $f'(0) = 100$.
Avec $g(x) = \sqrt{x}$ en $4$~: $g'(x) = \frac{1}{2\sqrt{x}}$, donc la
limite est $g'(4) = \frac14$.
\end{solution}

\begin{solution}{exo:g12:deriv:4}
\omterm{def:g11:func:function}{Ensemble de définition} $\R \setminus \{1\}$. Limites~: en $\pm\infty$,
$f(x) \sim x \to \pm\infty$~; en $1^\pm$, le numérateur $\to 4 > 0$ et le
dénominateur $\to 0^\pm$, donc $f \to \pm\infty$~; la droite $x = 1$ est
une \omterm{def:g12:limcont:asymptote}{asymptote verticale}.

\[
f'(x) = \frac{2x(x-1) - (x^2+3)}{(x-1)^2} = \frac{x^2 - 2x - 3}{(x-1)^2}
= \frac{(x-3)(x+1)}{(x-1)^2}.
\]
D'où $f' > 0$ sur $\intoo{-\infty}{-1}$ et $\intoo{3}{+\infty}$,
$f' < 0$ sur $\intoo{-1}{1}$ et $\intoo{1}{3}$. La \omterm{def:g11:func:function}{fonction} croît jusqu'à
un \omterm{def:g12:deriv:extremum}{maximum local} $f(-1) = -2$, décroît vers $-\infty$, saute vers
$+\infty$ après l'\omterm{def:g12:limcont:asymptote}{asymptote}, décroît jusqu'à un \omterm{def:g10:functions:extrema}{minimum} local
$f(3) = 6$, puis croît.
\end{solution}

\begin{solution}{exo:g12:deriv:5}
Pour $x \in \intoo{0}{15}$, la boîte a une base carrée de côté
$30 - 2x$ et une hauteur $x$, donc
\[
V(x) = x(30 - 2x)^2 .
\]
$V'(x) = (30-2x)^2 + x \cdot 2(30-2x)(-2) = (30-2x)(30 - 2x - 4x)
= (30-2x)(30-6x)$. Sur $\intoo{0}{15}$, $30 - 2x > 0$, donc $V'$ a le
signe de $30 - 6x$~: positif avant $x = 5$, négatif après. Le volume est
maximal pour $x = 5$~cm, donnant $V(5) = 5 \times 20^2 = 2000\
\text{cm}^3$.
\end{solution}

\begin{solution}{exo:g12:deriv:6}
\emph{1.} $f'(x) = 4x^3 - 12x^2$ et $f''(x) = 12x^2 - 24x = 12x(x-2)$.
Ainsi $f'' > 0$ sur $\intoo{-\infty}{0}$ et sur $\intoo{2}{+\infty}$
(\omterm{def:g12:deriv:convex}{convexe}), $f'' < 0$ sur $\intoo{0}{2}$ (\omterm{def:g12:deriv:convex}{concave}). $f''$ change de signe
en $0$ et en $2$~: ce sont tous deux des \omterm{def:g12:deriv:inflection}{points d'inflexion}.

\emph{2.} En $a = 2$~: $f(2) = 16 - 32 + 10 = -6$, $f'(2) = 32 - 48 =
-16$, \omterm{def:g12:deriv:tangent}{tangente} $y = -6 - 16(x-2) = -16x + 26$. L'écart est
\[
g(x) = f(x) + 16x - 26 = x^4 - 4x^3 + 16x - 16 .
\]
Comme $g(2) = g'(2) = g''(2) = 0$, $(x-2)^3$ divise $g$~; la division
donne $g(x) = (x-2)^3(x+2)$. Au voisinage de $x = 2$, $x + 2 > 0$, donc
$g$ a le signe de $(x-2)^3$~: négatif avant $2$, positif après. La courbe
traverse la \omterm{def:g12:deriv:tangent}{tangente}~: $2$ est bien un \omterm{def:g12:deriv:inflection}{point d'inflexion}.
\end{solution}

\begin{solution}{exo:g12:deriv:7}
La \omterm{def:g11:func:function}{fonction} $f(x) = \sqrt{x}$ satisfait
$f''(x) = -\frac{1}{4}x^{-3/2} < 0$ sur $\intoo{0}{+\infty}$~: elle est
\omterm{def:g12:deriv:convex}{concave}, donc son \omterm{def:g10:functions:graph}{graphique} est \emph{sous} chaque \omterm{def:g12:deriv:tangent}{tangente}. La \omterm{def:g12:deriv:tangent}{tangente}
en $a = 1$ est
\[
y = f(1) + f'(1)(x - 1) = 1 + \tfrac12 (x-1) = \frac{x+1}{2},
\]
d'où $\sqrt{x} \leq \frac{x+1}{2}$ pour tout $x > 0$. L'égalité signifie
que la courbe touche la \omterm{def:g12:deriv:tangent}{tangente}, ce qui n'arrive qu'au point de
contact $x = 1$. (Équivalent~: $\left(\sqrt x - 1\right)^2 \geq 0$.)
\end{solution}

\begin{solution}{exo:g12:deriv:8}
Par le \cref{thm:g12:deriv:convexity}, le \omterm{def:g10:functions:graph}{graphique} d'une \omterm{def:g11:func:function}{fonction}
\omterm{def:g12:deriv:convex}{convexe} \omterm{def:g12:deriv:derivative}{dérivable} est au-dessus de chacune de ses \omterm{def:g12:deriv:tangent}{tangentes}. La \omterm{def:g12:deriv:tangent}{tangente}
en $a$ est horizontale ($f'(a) = 0$), d'\omterm{def:g10:algebra:equation}{équation} $y = f(a)$. D'où
$f(x) \geq f(a)$ pour tout $x \in \R$~: $f(a)$ est le \omterm{def:g10:functions:extrema}{minimum} global.
\end{solution}

\begin{solution}{exo:g12:deriv:9}
\emph{Périmètre fixe.} Un rectangle de côtés $x$ et $\frac{p}{2} - x$
($0 < x < \frac p2$) a pour aire $S(x) = x\left(\frac p2 - x\right)$.
Alors $S'(x) = \frac p2 - 2x$ s'annule en $x = \frac p4$, positif avant
et négatif après~: le \omterm{def:g10:functions:extrema}{maximum} est en $x = \frac{p}{4}$, où les deux
côtés valent $\frac p4$ --- un carré.

\emph{Aire fixe.} Les côtés $x$ et $\frac Ax$ donnent le périmètre
$P(x) = 2\left(x + \frac Ax\right)$, avec
$P'(x) = 2\left(1 - \frac{A}{x^2}\right)$, négatif pour $x < \sqrt A$ et
positif après~: \omterm{def:g10:functions:extrema}{minimum} en $x = \sqrt{A}$, un carré encore.

\emph{Lien.} Les deux énoncés sont duaux. Supposons qu'un rectangle non
carré minimise le périmètre à aire fixe $A$~; le carré de même périmètre
aurait une aire $> A$ par le premier énoncé, et le rétrécir
homothétiquement à l'aire $A$ diminuerait strictement son périmètre ---
contredisant la minimalité. Chaque énoncé implique donc l'autre.
\end{solution}

\begin{solution}{pb:g12:deriv:1}
\textbf{1.} $30x\left(3x^2 + 1\right)^4$~; \quad
$\dfrac{x}{\sqrt{x^2 + 9}}$~; \quad
$-\dfrac{2x}{\left(x^2 + 1\right)^2}$.

\textbf{2.} $y(4) = 5$, \omterm{def:g10:reffunc:affine}{coefficient directeur} $\frac45$~:
\omterm{def:g12:deriv:tangent}{tangente} $y = 5 + \frac45(x - 4)$, c'est-à-dire
$y = \frac45 x + \frac95$.

\textbf{3.} $f'(x) = \frac{(x^2 + 1) - x \cdot 2x}{(x^2+1)^2}
= \frac{1 - x^2}{(x^2 + 1)^2}$~: négative, puis positive, puis
négative de part et d'autre de $-1$ et $1$~: \omterm{def:g10:functions:extrema}{minimum}
$f(-1) = -\frac12$, \omterm{def:g10:functions:extrema}{maximum} $f(1) = \frac12$, avec l'\omterm{def:g12:limcont:asymptote}{asymptote
horizontale} $0$ des deux côtés.

\textbf{4.} $g'(x) = 3x^2 - 6x = 3x(x - 2)$~: \omterm{def:g12:seq:monotonic}{croissante},
\omterm{def:g10:functions:variations}{décroissante}, \omterm{def:g12:seq:monotonic}{croissante} de part et d'autre de $0$ et $2$~; \omterm{def:g12:deriv:extremum}{maximum
local} $g(0) = 4$, \omterm{def:g10:functions:extrema}{minimum} local $g(2) = 0$. $g''(x) = 6x - 6$~:
\omterm{def:g12:deriv:convex}{concave} avant $1$, \omterm{def:g12:deriv:convex}{convexe} après~: inflexion en $(1, 2)$.

\textbf{5.} \omterm{def:g12:deriv:tangent}{Tangente} en $1$~: \omterm{def:g10:reffunc:affine}{coefficient directeur}
$g'(1) = -3$, d'où $y = 2 - 3(x - 1) = -3x + 5$. Différence~:
$x^3 - 3x^2 + 4 - (-3x + 5) = x^3 - 3x^2 + 3x - 1 = (x - 1)^3$,
qui change de signe en $1$~: la \omterm{def:g12:deriv:tangent}{tangente} \emph{traverse} la courbe
--- le comportement caractéristique d'un \omterm{def:g12:deriv:inflection}{point d'inflexion}, où la
courbe change de côté.

\textbf{6.} $T(x) = \dfrac{\sqrt{x^2 + 900}}{5} +
\dfrac{\sqrt{(40 - x)^2 + 400}}{2}$. Entrer avant $0$ ou au-delà de
$40$ rallonge les \emph{deux} portions~: c'est sans intérêt.

\textbf{7.} $T'(x) = \dfrac{x}{5\sqrt{x^2 + 900}} -
\dfrac{40 - x}{2\sqrt{(40 - x)^2 + 400}}$.

\textbf{8.} Dans le triangle de la course, le côté le long du
rivage vaut $x$ et l'hypoténuse $\sqrt{x^2 + 900}$~: le \omterm{def:g11:trigo:cossin}{sinus} de
l'angle avec la perpendiculaire est précisément leur quotient (côté
opposé sur hypoténuse)~; de même
$\sin\theta_2 = \frac{40 - x}{\sqrt{(40-x)^2 + 400}}$. Alors
$T'(x) = 0$ s'écrit
$\frac{\sin\theta_1}{5} = \frac{\sin\theta_2}{2}$ --- la \omterm{def:g11:prob:rv}{loi} de la
réfraction, avec les vitesses $5$ et $2$.

\textbf{9.} $T(24) \approx 20.5$ s~; $T(40) = 10 + 10 = 20$ s~;
$T(0) = 6 + \sqrt{2000}/2 \approx 28.4$ s. La ligne droite perd
même face à «~courir jusqu'au bout, puis nager droit devant~» ---
et les deux perdent face au trajet réfracté.

\textbf{10.} La \omterm{thm:g12:limcont:ivt}{dichotomie} sur $T'$ (négative en $24$, positive en
$40$) \omterm{def:g12:seq:limit}{converge} vers $x^* \approx 34$ m, avec
$T(x^*) \approx 19.5$ s~: environ une seconde de moins que la ligne
droite --- la différence entre un sauvetage et un drame.

\textbf{11.} Chaque durée partielle est une \omterm{def:g11:func:function}{fonction} \omterm{def:g12:deriv:convex}{convexe} de $x$
(une branche de $\sqrt{\text{trinôme}}$), donc $T$ est \omterm{def:g12:deriv:convex}{convexe}, et
un point critique d'une \omterm{def:g11:func:function}{fonction} \omterm{def:g12:deriv:convex}{convexe} \omterm{def:g12:deriv:derivative}{dérivable} en est le \omterm{def:g10:functions:extrema}{minimum}
global (\cref{exo:g12:deriv:8})~: $x^*$ n'est pas seulement
stationnaire, il est imbattable.

\textbf{12.} La lumière est plus lente dans l'eau~; d'après le
principe de Fermat, le trajet le plus rapide de l'air vers l'eau se
brise à la surface exactement selon la \omterm{def:g11:prob:rv}{loi} de la réfraction --- si
bien que les rayons issus d'un crayon immergé parviennent brisés à
l'œil, et le cerveau, qui prolonge les lignes droites, voit le
crayon cassé au niveau de la surface.

\textbf{13.} $(x^2)'' = 2 > 0$~: la \omterm{def:g11:func:function}{fonction} est \omterm{def:g12:deriv:convex}{convexe}. \omterm{def:g12:deriv:convex}{Convexité}
appliquée au \omterm{prop:g10:coordgeom:midpoint}{milieu}~: $f\left(\frac{a+b}{2}\right) \leq
\frac{f(a) + f(b)}{2}$, ce qui est exactement l'énoncé. À la main~:
$\frac{a^2 + b^2}{2} - \left(\frac{a+b}{2}\right)^2 =
\frac{(a - b)^2}{4} \geq 0$.

\textbf{14.} En prenant les racines carrées (\omterm{def:g11:func:monotone}{fonction croissante})
dans la question 13~:
$\frac{a + b}{2} \leq \sqrt{\frac{a^2 + b^2}{2}}$. Chaîne complète
sur $2$ et $8$~: harmonique $\frac{2 \cdot 16}{10} = 3.2$~;
géométrique $\sqrt{16} = 4$~; arithmétique $5$~; quadratique
$\sqrt{34} \approx 5.83$~: chaque \omterm{def:g11:stat:mean}{moyenne} s'incline devant la
suivante.

\textbf{15.} \omterm{def:g12:deriv:tangent}{Tangente} à $\frac1x$ en $1$~: valeur $1$, \omterm{def:g10:reffunc:affine}{coefficient
directeur} $-1$, d'où $y = 2 - x$. Les courbes \omterm{def:g12:deriv:convex}{convexes} se tiennent
au-dessus de leurs \omterm{def:g12:deriv:tangent}{tangentes}
(\cref{thm:g12:deriv:convexity})~: $\frac1x \geq 2 - x$ pour tout
$x > 0$, avec égalité exactement au point de contact $x = 1$.

\textbf{16.} À l'inflexion, le nombre \emph{quotidien} de cas (la
\omterm{def:g12:deriv:derivative}{dérivée}) atteint son \omterm{def:g10:functions:extrema}{maximum}~: la croissance cesse d'accélérer et
commence à ralentir --- le premier signal mathématique que la vague
se retourne, bien avant que les comptages eux-mêmes ne baissent.
Avant ce point, $f'' > 0$ (chaque jour pire que le précédent)~;
après, $f'' < 0$ (encore en croissance, mais de plus en plus
lentement).

\textbf{17.} $S'(r) = 4\pi r - \frac{2V}{r^2} = 0$ donne
$r^3 = \frac{V}{2\pi} = \frac{330}{2\pi}$~: $r \approx 3.74$ cm, et
$h = \frac{330}{\pi r^2} \approx 7.49$ cm.

\textbf{18.} $h = \frac{V}{\pi r^2} = \frac{2\pi r^3}{\pi r^2} = 2r$
à l'optimum~: la hauteur égale le diamètre, quel que soit le
volume.

\textbf{19.} $S''(r) = 4\pi + \frac{4V}{r^3} > 0$~: la \omterm{def:g11:func:function}{fonction} est
\omterm{def:g12:deriv:convex}{convexe}, donc le rayon critique donne le \omterm{def:g10:functions:extrema}{minimum} global. Les
canettes réelles sont plus élancées pour la prise en main, pour
l'empilement, pour la présence en rayon et à cause du métal plus
épais du dessus et du dessous --- l'optimisation minimise toujours
exactement ce que l'on a écrit, jamais ce que l'on voulait dire.

\textbf{20.} Maître nageur~: modéliser $T(x)$, dériver, réfraction,
\omterm{def:g12:deriv:convex}{convexité}, une seconde gagnée. Boîte~: modéliser $S(r)$, dériver,
$h = 2r$, \omterm{def:g12:deriv:convex}{convexité}, du métal économisé. Poutre
(\cref{pb:g11:deriv:1})~: modéliser $S(w)$, dériver,
$d = w\sqrt2$, tableau, de la rigidité gagnée. Et d'après le
principe de Fermat, la lumière elle-même déroule cette chaîne à
chaque surface qu'elle rencontre --- la nature fut le premier
optimiseur.
\end{solution}
